% Maintained model specification for the bounded sensing model.
% Exactness label: paper-exact.

To reflect the finite resolution and saturation of practical onboard sensors,
we define $d_{i,k}$ as the shortest Euclidean distance from vehicle $i$ to the
rectangular RoI of task $k$, regularized as
$\bar d_{i,k}=\max(d_{i,k},d_{\min})$. The bounded spatial sensing potential is

\begin{equation}
\bar A_{i,k}=\min\!\left\{
\delta_i(1-\mathcal O_{i,k})
\left(\frac{\bar d_{i,k}}{r_0}\right)^{-\zeta},
A_{\max}
\right\}.
\end{equation}

Unlike an unbounded inverse-distance model, $\bar A_{i,k}$ does not imply
arbitrarily large information gain when a vehicle is located close to or
inside an RoI. The corresponding sensing quality is

\begin{equation}
q_{i,k}(\tau_{i,k}^{s})=\bar A_{i,k}
\ln\!\left(1+\gamma s_i\tau_{i,k}^{s}\right).
\end{equation}

For candidate selection, we first count the number of coverage segments and
then apply the diversity factor:

\begin{equation}
S_k=\max\!\left(1,\left\lceil\frac{a_k}{2r_0}\right\rceil\right),\qquad
N_{\max}^{(k)}=\max\!\left(2,\left\lceil\varrho S_k\right\rceil\right).
\end{equation}

Candidates are ranked lexicographically by
$(\bar A_{i,k},A^{\rm raw}_{i,k},-i)$ in descending order. Thus the bounded
potential used by the sensing-quality model is the primary selection score;
the raw pre-saturation potential is consulted only to resolve ties at
$A_{\max}$, and the vehicle index provides a final deterministic tie-break.

Let $c_{i,k}\in\{0,1\}$ denote whether vehicle $i$ remains a candidate for
task $k$, and let $m_i=\max_k c_{i,k}$ be its vehicle-level eligibility
indicator. Candidate selection precedes the leader's bandwidth allocation.
The actor emits a bandwidth-utilization logit and vehicle logits $z_i$; the
feasible allocation is decoded as

\begin{equation}
w_i=
\begin{cases}
u\dfrac{m_i\exp(z_i)}{\sum_jm_j\exp(z_j)}, & \sum_jm_j>0,\\[4pt]
0, & \sum_jm_j=0,
\end{cases}
\qquad u=\operatorname{sigmoid}(z_u).
\end{equation}

Thus a vehicle excluded from every task receives no uplink bandwidth. When
candidate pruning is disabled, $m_i=1$ for every vehicle and the decoder
reduces to the unmasked sigmoid/softmax allocation.

The formal experiments use $d_{\min}=1$ m, $A_{\max}=50$, $r_0=80$ m,
$\varrho=4$, and $p_{\max}=0.3$. For the maintained short-RoI geometry, this
retains up to four ranked candidates per task, providing headroom for multiple
valid contributors while still pruning most vehicle--task pairs at the default
density. The reward-penalty weights are
$(\omega_{\mathrm{fail}},\omega_{\mathrm{lat}},
\omega_{\mathrm{part}})=(2,3,3)$. These parameters are applied identically to
all applicable learning-based methods, together with the fixed common learning
rate, clipping range, training budget, and network widths declared in the
experiment protocol.

Task completion is evaluated by
$I_k=\mathbb{I}(Q_k>0)$ and is reported both overall and separately for the
high-, medium-, and low-urgency classes. Safety-critical and traffic-efficiency
tasks are treated as protected services. A low-urgency environment-maintenance
task, however, may remain unserved when its marginal sensing benefit does not
justify the required payment, energy consumption, and bandwidth cost. This is
a permitted priority-aware allocation outcome rather than a feasibility
violation; it is still counted as a non-completed task, and the corresponding
low-urgency completion rate is disclosed explicitly.

% Resource and mechanism analyses.
Resource-sensitivity results use selected checkpoints. Unless otherwise stated, each point first averages over
32 independent scheduling-cycle realizations and is then summarized over eight
independent training seeds; error bars denote 95\% confidence intervals across
training seeds. For the CPU-capacity study, every learning-based method is
trained once under the same episode-level domain randomization: one of five
CPU-capacity intervals centered at 5, 10, 15, 20, and 25 Gcycles/s is sampled
uniformly for each scheduling cycle, after which the vehicles' $F_i$ values are
sampled independently within that interval. The selected policy is then tested
at each fixed interval; no method is retrained separately for an individual
CPU-capacity point. The three panels report social welfare, RSU utility, and
task-completion rate, respectively.
The bandwidth-robustness panel reuses the bandwidth sweep, which uses
five common training seeds and reports its confidence intervals across those
five seed-level means.

For pricing sensitivity, we hold the learned bandwidth allocation fixed and
perturb only the learned price according to
\begin{equation}
p_k(\vartheta)=\min\!\left\{p_{\max},
\max\!\left(0,\vartheta p_k^{\mathrm{TE}}\right)\right\}.
\end{equation}
The analytical follower responses are recomputed after each perturbation.
Source data additionally report the fraction of prices clipped by $p_{\max}$.

The compression-energy panel is a component-level counterfactual analysis.
For vehicles selected by the unperturbed TE-PPO policy, the sampled vehicle and
channel state, bandwidth allocation, and sensing duration are held fixed while
only the compression ratio is varied. Consequently, the panel isolates the
opposing trends of local computation energy and uplink transmission energy; it
does not claim that every counterfactual compression point is a feasible or
re-optimized policy response.
